\chapter{Decodifica vincolata: il componente simbolico}
\label{cap:decoding}

Questo capitolo descrive l'unico elemento propriamente \emph{neuro-simbolico} del
sistema, e riporta il risultato più controintuitivo dell'intero lavoro: il
vincolo simbolico \emph{peggiora} la metrica di sovrapposizione e
contemporaneamente \emph{abilita} il segnale di reward.

\section{Il problema: dal vincolo sulle parole al vincolo sui token}

L'uscita valida è una sequenza di gloss appartenenti a un vocabolario finito $G$
di $15\,472$ elementi. Il modello, però, genera \emph{token}, e un gloss
corrisponde in generale a più token (sezione sui fondamenti di tokenizzazione).
Non si può quindi semplicemente restringere il vocabolario di uscita: bisogna
tradurre il vincolo sul livello delle parole in un vincolo sul livello dei token.

La formulazione corretta è quella della~\eqref{eq:trie-allowed}: al passo $t$,
dato il prefisso di token già emesso, sono ammessi i token che mantengono il
prefisso estendibile a un elemento di $G$. Il Trie costruito su $G$ è la
struttura dati che rende questo calcolo immediato.

Operativamente il vincolo agisce sui logits come nella~\eqref{eq:constrained}:
si pone $z_{t,v} = -\infty$ per ogni $v \notin \allowed_t$ prima della softmax.

\section{Implementazione}

Due componenti, con responsabilità separate:

\begin{center}
\begin{tabular}{ll}
\toprule
File & Ruolo \\
\midrule
\texttt{src/grammar/gloss\_grammar.py} & \texttt{GlossVocabularyMask}: il Trie e il calcolo dell'insieme ammesso \\
\texttt{src/grammar/grammar\_logits\_processor.py} & \texttt{GlossVocabularyLogitsProcessor}: l'applicazione ai logits \\
\bottomrule
\end{tabular}
\end{center}

La separazione è deliberata: la maschera è una struttura dati pura, testabile
senza modello, mentre il processore di logits è l'adattatore verso l'interfaccia
di generazione della libreria. Tutti i test sulla correttezza del vincolo agiscono
sulla prima, senza caricare la GPU.

\subsection{La doppia radice}
\label{sec:trie-dual-root}

Un dettaglio che non è un dettaglio. Il tokenizzatore codifica lo spazio come
parte del token successivo: il token per \texttt{"HOUSE"} e quello per
\texttt{"~HOUSE"} sono identificatori distinti. Ne segue che il primo gloss di
una sequenza e i gloss successivi hanno insiemi di primi token \emph{diversi}.

Il Trie è quindi costruito con \textbf{due radici}:
\begin{description}
  \item[\texttt{root}] usata per il primo gloss della sequenza, i cui token
        iniziali non hanno spazio davanti;
  \item[\texttt{space\_root}] usata per tutti i gloss successivi, i cui token
        iniziali portano lo spazio.
\end{description}

Se si usasse una sola radice si aprirebbe un difetto concreto: il modello
potrebbe concatenare due gloss senza separatore, producendo forme come
\texttt{DEBUTRECHT} (dalla giustapposizione di \texttt{DEBUT} e
\texttt{RECHT}) che il Trie a radice singola accetterebbe come prefisso valido
perché non distingue il confine di parola. La doppia radice rende il confine
esplicito nella struttura.

\subsection{Numeri della costruzione}

\begin{center}
\begin{tabular}{lr}
\toprule
Grandezza & Valore \\
\midrule
Gloss nel vocabolario & $15\,472$ \\
Identificatori di token ammessi come primo token & $3\,967$ \\
Gloss con il primo token bloccato & $4$ \\
Massa di probabilità irraggiungibile & $12 / 897\,971$ \\
\bottomrule
\end{tabular}
\end{center}

I quattro gloss non emettibili sono \texttt{-RRB-}, \texttt{WWW.8HOURS.EU},
una forma costituita da un asterisco isolato, e una forma che inizia con il
carattere di \emph{byte order mark}. Sono tutti artefatti del corpus
(sezione~\ref{sec:artefatti}), non gloss linguisticamente significativi: la loro
inaccessibilità non costituisce una perdita.

La massa irraggiungibile di $12$ su $897\,971$ è la verifica quantitativa che il
vincolo non stia tagliando fuori qualcosa di rilevante: circa una parte su
settantacinquemila.

\subsection{L'interfaccia per gli usi non generativi}

Il vincolo serve anche al di fuori della generazione: gli obiettivi ausiliari
del Capitolo~\ref{cap:ausiliari} hanno bisogno di conoscere l'insieme ammesso in
modalità \emph{teacher-forced}, cioè per prefissi noti in blocco e non generati
uno alla volta. A questo scopo è stato aggiunto il metodo
\begin{center}
\texttt{allowed\_mask\_for\_prefixes(prefixes, vocab\_size, device)}
\end{center}
con il supporto interno \texttt{\_allowed\_after\_prefix}. Restituisce
direttamente la maschera booleana per un insieme di prefissi, evitando di
duplicare la logica del Trie in due punti del codice: la definizione di ``token
ammesso'' esiste in un solo posto, e sia la generazione sia la loss ausiliaria
la consumano.

\section{L'alternativa scartata: l'automa a stack}
\label{sec:pda-scartato}

Il progetto conteneva, in una fase precedente, un'implementazione di automa a
stack (\emph{pushdown automaton}) con parsing LL(1), pensata per imporre una
grammatica context-free invece di un semplice lessico chiuso. È stata rimossa.
Le ragioni vanno documentate, perché la rimozione di codice funzionante richiede
giustificazione.

\begin{enumerate}
  \item \textbf{Non è mai stata eseguita.} Il componente non ha mai prodotto un
        risultato sperimentale. Manteneva quindi un costo di manutenzione senza
        contropartita.
  \item \textbf{Conflitto di parsing non risolto.} La grammatica presentava un
        conflitto fra la produzione vuota e l'insieme dei simboli che possono
        seguire un non terminale, innescato dai gloss che iniziano con una cifra
        (per esempio \texttt{0}, \texttt{8}, \texttt{08.30}). Un parser LL(1)
        richiede che l'insieme dei primi simboli e quello dei simboli seguenti
        siano disgiunti in presenza di produzioni vuote; qui non lo erano.
  \item \textbf{Sovradimensionamento.} Il linguaggio target è
        \texttt{gloss*}: una chiusura di Kleene su un insieme finito. È un
        linguaggio regolare, riconoscibile da un automa a stati finiti. Un automa
        a stack riconosce linguaggi context-free, strettamente più espressivi:
        la potenza aggiuntiva non serviva a nulla su questo compito.
\end{enumerate}

La conclusione generale è che la scelta del formalismo simbolico deve seguire la
complessità effettiva del linguaggio da vincolare, non l'ambizione del disegno.
Sulla decodifica grammaticalmente vincolata in senso pieno esiste letteratura
matura~\cite{geng2023gcd,park2024gad}; sarebbe la strada corretta se il vincolo
fosse davvero context-free, che qui non è il caso.

\section{Il risultato controintuitivo}
\label{sec:trie-risultato}

Ecco l'effetto del vincolo sulla cella base zero-shot.

\begin{center}
\begin{tabular}{lrr}
\toprule
Metrica & Senza Trie & Con Trie \\
\midrule
ROUGE-L & $0{,}3648$ & $0{,}1373$ \\
Non-copy-token accuracy & $0{,}0006$ & $0{,}0432$ \\
Validity & $0{,}0370$ & $0{,}9055$ \\
\bottomrule
\end{tabular}
\end{center}

Il vincolo fa \emph{crollare} ROUGE-L da $0{,}3648$ a $0{,}1373$: una perdita
apparente di $0{,}23$. Allo stesso tempo la non-copy-token accuracy passa da
$0{,}0006$ a $0{,}0432$, un fattore circa $70$.

\subsection{Come si spiega}

Le due metriche misurano cose diverse, e il vincolo agisce su di esse in
direzioni opposte.

Senza vincolo, il modello base fa la cosa per cui è stato pre-addestrato: produce
inglese. Poiché il gloss è in larga parte inglese maiuscolizzato
(sezione~\ref{sec:sintetico}), produrre inglese ottiene un ROUGE-L
rispettabile ``per caso''. Ma la non-copy-token accuracy, che valuta solo le
posizioni dove il gloss \emph{differisce} dal testo, è $0{,}0006$: praticamente
nulla. Il modello non ha imparato niente sul gloss; sta beneficiando della
sovrapposizione lessicale del corpus.

Con il vincolo, il modello è costretto a emettere solo gloss del vocabolario.
Perde la possibilità di copiare l'inglese, quindi perde il ROUGE-L gratuito. In
compenso, sulle posizioni che richiedono una trasformazione reale, comincia a
produrre qualcosa di corretto ($0{,}0432$ invece di $0{,}0006$).

Detto in modo compatto: il vincolo trasforma una copia con punteggio alto in un
tentativo di traduzione con punteggio basso. La metrica di sovrapposizione
registra un peggioramento; la metrica che conta registra un miglioramento di
quasi due ordini di grandezza.

\subsection{L'effetto veramente importante: il vincolo crea il supporto}

Il ROUGE-L, però, non è la ragione per cui il vincolo è presente nel sistema. La
ragione è nella tabella del supporto dei rollout
(sezione~\ref{sec:supporto}), che vale la pena riportare nella parte pertinente:

\begin{center}
\begin{tabular}{lrr}
\toprule
Substrato zero-shot & Gruppi morti & Reward media \\
\midrule
Senza Trie & $0{,}7390$ & $-0{,}9579$ \\
Con Trie   & $0{,}1990$ & $-0{,}8034$ \\
\bottomrule
\end{tabular}
\end{center}

Senza vincolo, il $74\%$ dei gruppi è morto: come stabilito nella
sezione~\ref{sec:gruppi-morti}, per quei prompt il gradiente è esattamente zero e
il reinforcement learning non è nemmeno definito come procedura utile. Con il
vincolo, i gruppi morti scendono al $20\%$.

\textbf{Il componente simbolico non migliora le metriche: rende possibile
l'apprendimento.} È il risultato neuro-simbolico di questo lavoro, e ha una forma
argomentativa precisa. Il vincolo restringe il supporto della distribuzione
generativa alla regione in cui la reward è \emph{discriminante}. Fuori da quella
regione tutti i candidati sono ugualmente cattivi, quindi la
sottrazione~\eqref{eq:grpo-adv} li annulla tutti; dentro, differiscono, quindi
producono gradiente.

Questo giustifica anche una raccomandazione operativa generale: quando si prepara
un esperimento di reinforcement learning su un compito con uscita strutturata,
conviene misurare il supporto dei rollout \emph{prima} di addestrare, e
considerare il vincolo simbolico non come un accessorio di qualità ma come una
precondizione per la fattibilità.

\section{Il costo}

Il vincolo non è gratuito. A ogni passo di generazione richiede:
\begin{itemize}
  \item la risoluzione del nodo corrente del Trie a partire dal prefisso;
  \item la costruzione di una maschera booleana sulla dimensione del vocabolario;
  \item l'applicazione della maschera ai logits.
\end{itemize}

Il costo dominante è la costruzione della maschera, lineare in $|\vocab|$. Nel
regime di questo progetto (modello piccolo, $G=8$ candidati per gruppo,
completion di al massimo $128$ token) resta trascurabile rispetto al costo del
passaggio in avanti della rete. Su modelli più grandi la valutazione andrebbe
rifatta, e la memorizzazione delle maschere per nodo del Trie sarebbe
l'ottimizzazione naturale.

\section{Che cosa non fa il vincolo}

Per evitare sovrainterpretazioni, va delimitato ciò che il Trie garantisce.

\begin{description}
  \item[Garantisce] che ogni parola emessa appartenga al vocabolario gloss e che
        i confini di parola siano rispettati.
  \item[Non garantisce] che la \emph{sequenza} sia grammaticalmente valida in
        ASL: il linguaggio riconosciuto è \texttt{gloss*}, cioè qualunque
        successione di gloss legali.
  \item[Non garantisce] che la sequenza sia una traduzione corretta della frase
        di ingresso. Il vincolo è sintattico e non guarda l'ingresso.
  \item[Non impedisce] le ripetizioni degeneri: \texttt{HOUSE HOUSE HOUSE} è una
        sequenza perfettamente ammissibile per il Trie. È il motivo per cui la
        reward include una componente dedicata alla ripetizione
        (Capitolo~\ref{cap:reward}) e per cui il test avversario sul peso $w$
        della sezione precedente era necessario.
\end{description}

Il vincolo è dunque una condizione necessaria di buona formazione, non una
condizione sufficiente di correttezza. Confondere le due cose porterebbe a
attribuire al componente simbolico garanzie che non ha.

\section{Il costo nascosto della generazione vincolata}

Il risultato della sezione~\ref{sec:trie-risultato} --- il vincolo peggiora la
metrica di sovrapposizione --- non è un'anomalia locale di questo progetto. È
un fenomeno riconosciuto nella letteratura sulla generazione strutturata, e
conviene collocarvelo perché ne chiarisce la natura.

Il punto teorico è quello anticipato a proposito della~\eqref{eq:constrained}:
mascherare i logits e rinormalizzare produce una distribuzione che \emph{non} è
la distribuzione condizionata del modello all'insieme delle sequenze valide. La
rinormalizzazione è locale, passo per passo, mentre il condizionamento corretto
richiederebbe di tenere conto della probabilità che il prefisso corrente possa
essere completato in una sequenza valida. Ne segue un errore sistematico, ed
esistono metodi che lo correggono restituendo stimatori asintoticamente non
distorti~\cite{ye2504}.

Le conseguenze pratiche di questa distorsione sono state documentate da due
angoli distinti. Da un lato, il vincolo può degradare la qualità del contenuto
generato anche quando garantisce la forma, con un costo che è stato
caratterizzato come dipendente dal condizionamento sulla bozza
\cite{reddy2603}. Dall'altro, il vincolo può indurre il modello a proseguire
lungo una struttura formalmente valida ma sostanzialmente sbagliata, un effetto
descritto come tassa di allineamento della decodifica
vincolata~\cite{zhou2604}.

Nel presente lavoro la manifestazione è precisamente questa: il modello base
vincolato produce sequenze di gloss ben formate (validity $0{,}9055$) ma con
punteggio di sovrapposizione basso ($0{,}1373$), mentre senza vincolo produce
sequenze mal formate (validity $0{,}0370$) con punteggio più alto ($0{,}3648$).
La differenza rispetto alla letteratura citata è nella conclusione operativa:
qui il vincolo va mantenuto \emph{nonostante} il costo, perché è ciò che rende
non nullo il gradiente del reinforcement learning
(sezione~\ref{sec:supporto}). Il costo sulla metrica di sovrapposizione è il
prezzo pagato per avere un segnale di apprendimento.

\section{Un secondo meccanismo simbolico: la riparazione post-hoc}
\label{sec:rule-repair}

Il Trie vincola \emph{durante} la generazione: sceglie fra le continuazioni
ammesse, ma non ha nozione della frase sorgente e non può preferire un gloss
corretto a uno semplicemente legale. Questa sezione descrive un secondo
meccanismo simbolico, complementare e indipendente, che agisce \emph{dopo} la
generazione e usa esattamente l'informazione che il Trie non ha: il testo
inglese di partenza.

\subsection{La motivazione: dove sbaglia ancora la policy}

Ispezionando le generazioni salvate della policy SFT (la più competente della
campagna) emerge un pattern netto: gli errori residui sono quasi tutti
sostituzioni lessicali puntuali su parole sorgente rare ---
\texttt{GUISE}$\to$\texttt{GUARD}, \texttt{BLINDNESS}$\to$\texttt{BLIGHT},
\texttt{DESC-LIVE}$\to$\texttt{LIVE}. I gloss d'oro che il modello manca sono
\emph{hapax legomena} del train set: frequenza mediana $1$, il $61\%$ presente
tre volte o meno, contro una frequenza di base dell'$1{,}6\%$ per i gloss
d'oro in generale.

La conseguenza è doppia. Un modello neurale non può memorizzare un hapax da un
adattamento LoRA su un corpus di questa taglia; e il reinforcement learning a
gradiente di policy non può correggerlo nemmeno in linea di principio, perché
sul $71{,}7\%$ dei fallimenti il token corretto non compare in nessuno dei
campioni generati per quel prompt: nessuna reward può assegnargli credito, per
la stessa ragione strutturale del collasso del supporto discusso nella
sezione~\ref{sec:supporto}.

Un transduttore \emph{deterministico}, però, non ha questo problema: mappa la
parola sorgente direttamente, indipendentemente dal fatto che sia mai stata
osservata durante l'addestramento. I due sistemi hanno quindi profili di
errore complementari, ed è questo che il meccanismo sfrutta.

\subsection{Il meccanismo}

Per ogni token generato dal modello si chiede se è \emph{derivabile} dalla
frase sorgente: condividerne lo stem (almeno $4$ caratteri di prefisso comune,
dopo aver rimosso i marcatori \texttt{DESC-}/\texttt{X-}/\texttt{fs-}) con
almeno una parola del testo. Se non lo è --- il modello ha probabilmente
allucinato un elemento lessicale --- viene sostituito con il token allineato
prodotto dal transduttore a regole del Capitolo~\ref{cap:dataset}; se lo è, il
token del modello resta intatto.

La sostituzione è ristretta alle posizioni allineate uno a uno fra l'uscita del
modello e quella del transduttore (via \texttt{difflib.SequenceMatcher}):
inserzioni e cancellazioni restano decisioni del modello, che gestisce la
lunghezza e le parole funzionali molto meglio del sistema a regole. La mappa
forma-sorgente$\to$forma-d'oro (\texttt{fit\_morphology}) sostituisce un
lemmatizzatore esterno, appresa dalla sola morfologia osservata nel train
(\texttt{hidden}$\to$\texttt{HIDE}, \texttt{outsourcing}$\to$\texttt{OUTSOURCE}):
oltre a evitare una dipendenza offline sul cluster, ottiene un exact match più
alto di WordNet sulla baseline a regole ($64{,}90\%$ contro $63{,}10\%$).
L'implementazione è in \texttt{src/analysis/rule\_repair.py}; lessico e
morfologia sono fittati esclusivamente sul train split.

\subsection{Effetto misurato}

Sulle generazioni salvate di cinque celle indipendenti, prima campione per
prompt, $2000$ prompt di eval:

\begin{center}
\begin{tabular}{lrrrrr}
\toprule
Celle & EM prima & EM dopo & Corrette & Rotte & $p$ (McNemar) \\
\midrule
SFT & $87{,}10\%$ & $88{,}85\%$ & $40$ & $5$ & $7{,}9\times 10^{-8}$ \\
SFT (passata few-shot) & $70{,}20\%$ & $75{,}60\%$ & $111$ & $3$ & $2{,}4\times 10^{-29}$ \\
GRPO few-shot & $5{,}20\%$ & $8{,}25\%$ & $61$ & $0$ & $8{,}7\times 10^{-19}$ \\
GRPO zero-shot (passata) & $1{,}85\%$ & $2{,}90\%$ & $21$ & $0$ & $9{,}5\times 10^{-7}$ \\
Ablation dr-grpo & $3{,}65\%$ & $5{,}55\%$ & $38$ & $0$ & $7{,}3\times 10^{-12}$ \\
\bottomrule
\end{tabular}
\end{center}

La metrica primaria non-copy-token accuracy si muove nella stessa direzione
(SFT $0{,}9793\to 0{,}9809$; GRPO few-shot $0{,}4642\to 0{,}5239$): non è uno
scambio fra una metrica e l'altra.

\subsection{Il caveat sulla generalizzazione, dichiarato}

La soglia di $4$ caratteri e l'euristica \texttt{is\_derivable} sono state
tarate ispezionando le uscite dell'SFT. La prova di generalizzazione è che lo
stesso meccanismo, non ritoccato, migliora anche i quattro insiemi su cui
\emph{non} è stato tarato (few-shot SFT, entrambe le celle GRPO, dr-grpo), e
non rompe mai più di $5$ prompt su $2000$ in nessun caso. Il tipo di
regressione è anch'esso compreso e circoscritto: un'abbreviazione la cui
espansione compare nel sorgente (\texttt{EU} da \emph{European Union}) non è
considerata derivabile secondo l'euristica dello stem, quindi il transduttore
la sovrascrive; è per questo che il rapporto corrette/rotte resta comunque
intorno a $10$:$1$ invece di essere privo di regressioni.

\subsection{Stato di integrazione}

Il meccanismo è cablato nella pipeline di valutazione come blocco opzionale e
\textbf{additivo}: la chiave \texttt{evaluation.rule\_repair} (default
\texttt{false}) applica la riparazione a ogni completion e ricalcola l'intero
blocco di metriche primarie sulle completion riparate, in un blocco separato
\texttt{results["rule\_repair"]} dei risultati JSON --- mai sovrascrivendo le
metriche primarie, sullo stesso modello del blocco \texttt{oracle\_best\_of\_n}
già esistente. La composizione è coperta da test unitari
(\texttt{tests/test\_rule\_repair\_eval\_wiring.py}); la tabella sopra resta
però una misura fatta sulle generazioni già salvate dalle run del cluster, non
ancora riprodotta end-to-end attraverso il nuovo flag su una valutazione dal
vivo. È il prossimo controllo naturale, a costo quasi nullo (nessun nuovo
addestramento, un solo passaggio di eval).

\section{Un terzo meccanismo: il glossario iniettato SOLO in addestramento}
\label{sec:glossario}

La sezione precedente ripara l'allucinazione lessicale \emph{dopo} la
generazione. Un'alternativa naturale è prevenirla, fornendo al modello
l'informazione lessicale mancante \emph{dentro} il prompt: un blocco
``Glossary:'' che elenca, per le parole rare della frase da tradurre, la
forma gloss corretta. Questa sezione discute due varianti non equivalenti, la
ragione per cui una è stata scartata sul nascere, e il disegno --- diverso da
entrambe quelle note in letteratura --- effettivamente implementato per
l'altra.

\subsection{Due varianti, non equivalenti}

\textbf{Iniezione a tempo di inferenza (zero-shot).} Aggiungere il blocco
glossario al prompt solo in fase di valutazione, senza toccare
l'addestramento, è l'opzione più economica da testare, ma ha un difetto di
fondo per un modello di questa taglia: la fase SFT di Qwen2.5-0.5B-Instruct su
questo compito non ha \emph{mai} visto un blocco ``Glossary:'' nel proprio
input (sezione~\ref{sec:sft-vs-regola-analisi}); \texttt{build\_t2g\_prompt}
produce, in modalità zero-shot, il testo grezzo della frase e nient'altro.
Testare un'iniezione di questo tipo in inferenza misurerebbe la robustezza del
modello a un formato di prompt mai osservato, non il potenziale reale
dell'idea: un risultato negativo sarebbe ambiguo fra ``il glossario non aiuta
questo compito'' e ``il modello non sa sfruttare una struttura di prompt che
non ha mai visto durante il fine-tuning'', e le due letture richiedono
conclusioni opposte. Per questo questa variante resta scartata: è lo stesso
principio metodologico del Capitolo~\ref{cap:ausiliari}, qualificare un
meccanismo prima di misurarlo.

\textbf{Annotazione a tempo di addestramento.} L'alternativa con precedenti
diretti in letteratura è addestrare il modello a \emph{usare}
l'annotazione, non solo a riceverla. Dinu et al.~\cite{dinu2019training}
mostrano che inserire vincoli terminologici come annotazioni inline
nell'input di addestramento --- non come post-processing né come vincolo di
decodifica --- permette a un NMT di imparare un comportamento di copia
affidabile verso il termine fornito, a costo di inferenza nullo rispetto alla
decodifica vincolata classica~\cite{hokamp2017,post2018}. Nella letteratura
più recente orientata ai modelli linguistici generici, l'iniezione puramente
a tempo di inferenza \emph{funziona}, ma su modelli e compiti diversi da
questo: DiPMT~\cite{ghazvininejad2023dipmt} ottiene guadagni fornendo
traduzioni candidate per le parole rare direttamente nel prompt di un LLM
istruito e multilingue, e Bogoychev e Chen~\cite{bogoychev2023terminology}
combinano un modello addestrato a rispettare la terminologia con un
raffinamento a base di prompting. In entrambi i casi il modello di partenza è
un sistema molto più grande e/o già esposto a task di traduzione con
terminologia a tempo di addestramento o pre-addestramento massivo: condizioni
che qui non valgono, il che motiva un disegno diverso invece di una copia
diretta di uno dei due.

\subsection{Il disegno implementato, e perché differisce da Dinu et al.}
\label{sec:glossario-disegno}

Il meccanismo (\texttt{src/utils/glossary.py}) inietta il blocco glossario
\textbf{solo nel prompt di addestramento GRPO}, mai in valutazione:
\texttt{eval\_t2g.py} non importa quel modulo, per costruzione, e ogni cella
resta valutata esattamente come le altre. Questo è deliberatamente diverso
dal meccanismo di Dinu et al., dove l'annotazione è presente sia in
addestramento sia in inferenza (un vero comportamento di copia): fornire il
glossario anche in eval userebbe il gold per selezionare l'aiuto, la stessa
fuga di informazione che l'esperimento vuole escludere, e testerebbe la
robustà a un formato di prompt piuttosto che il contributo
dell'addestramento. Il disegno qui misura quindi una domanda più stretta e
più onesta: \emph{l'esposizione a mappature corrette durante l'addestramento
migliora la gestione delle parole rare del modello DA SOLO, quando l'aiuto
sparisce}?

Due scelte tecniche derivano da questo vincolo:

\begin{description}
  \item[Origine del glossario: dal gold, non un embedding.] Le parole
        rare sono esattamente la popolazione che la
        sezione~\ref{sec:rule-repair} già caratterizza (hapax legomena, poca
        o nessuna occorrenza nel train): è precisamente dove un embedding
        semantico sarebbe meno affidabile, per la stessa scarsità di segnale
        che il glossario dovrebbe compensare. La soluzione adottata sfrutta
        un fatto specifico dell'addestramento GRPO: il gold di
        \emph{questo} esempio è già disponibile (è ciò su cui la reward
        valuta), quindi la forma corretta si legge direttamente dal gold
        della riga corrente
        (\texttt{align\_source\_to\_gold}, la stessa euristica di
        allineamento per stem di \texttt{rule\_repair.py}, applicata riga per
        riga invece che fittata sul corpus). Nessuna tabella pre-fittata,
        nessun embedding.
  \item[Dropout del blocco (\texttt{glossary.dropout}, default $0{,}5$).]
        Senza dropout, il modello vedrebbe il glossario in \emph{ogni}
        esempio di training che contiene parole rare e mai in valutazione:
        un disallineamento di distribuzione che rischierebbe di essere
        \emph{dannoso} (il modello impara a leggere un blocco che poi
        sparisce sempre). Omettendo il blocco su una frazione degli esempi
        --- pur avendone diritto --- la generazione senza glossario resta
        in-distribuzione anche durante l'addestramento stesso, e il confronto
        in eval (sempre senza glossario, per ogni cella) diventa
        interpretabile in entrambe le direzioni di un risultato.
\end{description}

Poiché la cella estende \texttt{grpo/zero-shot.yaml} o
\texttt{grpo/few-shot.yaml} (nessuna fase SFT), il rischio di riuso
silenzioso dell'impronta SFT sollevato in una stesura precedente di questa
sezione \textbf{non si applica}: non esiste un adattatore SFT da confondere
fra una cella con e senza glossario.

\subsection{Stato: implementato, non ancora addestrato}

Il meccanismo è verificato a livello di codice
(\texttt{tests/test\_glossary.py}, ed estensioni a
\texttt{tests/test\_rule\_repair.py} e \texttt{tests/test\_prompting.py} per
l'allineamento per riga e la composizione del prompt), con l'invariante
``blocco assente $\Rightarrow$ prompt byte-identico'' verificata esplicitamente,
sullo stesso modello del contratto di non regressione degli obiettivi
ausiliari (sezione~\ref{sec:contratto-ausiliari}). Due celle sono pronte
in \texttt{experiments/configs/qwen25-05b/ablations/glossary/} (zero-shot e
few-shot), fuori dalla campagna \texttt{--ablation} ufficiale a 15 celle per
non alterarla silenziosamente.

Il criterio di promozione, dichiarato in anticipo per coerenza col
Capitolo~\ref{cap:ausiliari}: la cella deve ridurre l'errore
\emph{specificamente} sulla popolazione di token gold hapax legomena
(la stessa che la sezione~\ref{sec:rule-repair} isola) più di quanto non
riduca l'errore genericamente sul resto del vocabolario. Se il guadagno (o
l'assenza di guadagno) è uniforme sulle due popolazioni, l'effetto --- se
presente --- è di regolarizzazione generica dell'aggiunta di più contesto al
prompt, non del meccanismo di glossario in sé. Nessuna delle due celle è
stata ancora addestrata mentre questa relazione viene scritta: il costo di
uno slot di cluster non è stato speso su un'idea la cui implementazione non
fosse già stata verificata a livello di codice.
